\chapter{卷积神经网络}

上一节，我们研究了线性神经网络，能够完美的解决MNIST-10分类问题了。那么，我们来看一看更大型的问题。
\begin{mdframed}[frametitle=例：ImageNet 10-分类问题]
ImageNet\cite{2009ImageNet}是一个大型的图像数据库，包含了超过1400万张标注的图像，分为1000个类别。这些图像的大小是不固定的，但我们都将其处理成$224\times 224 \times 3$的大小，并仅取其中的十类图像，将这些图像合并成一个数据集，称为ImageNet-10。

问：怎样解决这个分类问题？
\end{mdframed}

我们上一节的处理方式是：将输入二阶张量展平成一阶张量（\mintinline{python}|x.view(-1)|），然后输入到一个线性神经网络中。对于MNIST-10问题，这种方式是可行的，因为MNIST-10的图像大小为$28\times 28$，展平后是一个长度为784的一维向量，输入到一个线性神经网络中，参数量大约为$784 \times 10 = 7840$，这是可以接受的。

但这个问题中，这么处理就不可行了：$224\times 224 \times 3 = 150528$，展平后是一个长度为150528的一维向量，输入到一个单层线性分类器中，参数量大约为$150528 \times 10 = 1505280$。

一个线性层就需要1505280个参数！如果这样做，模型训练起来将会困难得多，我们需要一个新的手段来处理这些数据。

\section{卷积层}

\subsection{平移等变性和空间局部性}

面对图像数据，我们不禁想到了人眼。人眼在识别图像时，对于图像的不同地方采取的是相同的处理方式，即使这部分平移了，我们依然认识；不会像线性层那样，将某些特征和它们在图像中的绝对位置绑定在一起。我们希望神经网络也能像人眼一样，具有平移等变性，如图\ref{fig:conv-meow}所示。

另外，众所周知，很多图像是具有空间局部性的：某个像素点，和周围的像素点放在一起，能够表示一些深层的语义信息。比如，某个像素点和周围的像素点放在一起，可能表示一个眼睛；而如果这个像素点和周围的像素点放在一起，可能表示一个鼻子。我们希望神经网络能够利用这种空间局部性。

\begin{figure}
    \centering
    \includegraphics[width=0.3\textwidth]{conv-meow.jpg}
    \caption{平移等变性实例}
    \label{fig:conv-meow}
\end{figure}

这让我们想到了一个操作：如果能取一个“窗口”在图像上滑动，每次取窗口内的像素点进行相同的处理，输出一张新的“特征图像”，那么就可以实现平移等变性和空间局部性。用更严谨的数学语言来说，这个操作就是\textbf{用核函数进行局部加权映射}。这不禁让我们想到了数学中的卷积操作。

在数学分析中，我们学习过卷积公式：
\begin{equation}
    (f * g)(t) = \int_{-\infty}^{\infty} f(\tau) g(t - \tau) d\tau
\end{equation}
这又被称作是“一维连续卷积”。我们如果将其写成离散形式，应该是这样的：
\begin{equation}
    (f * g)(t) = \sum_{i=-\infty}^{\infty} f(i) g(t - i)
\end{equation}
把它推广到二维空间中去：
\begin{equation}
    (f * g)(x, y) = \sum_{i=-\infty}^{\infty} \sum_{j=-\infty}^{\infty} f(i, j) g(x - i, y - j)
\end{equation}
这就是\textbf{二维离散卷积公式}。

然而，严格遵照这个公式会带来一个小小的“麻烦”：注意求和项中的 $g(x-i, y-j)$。如果我们把 $g$ 视为一个尺寸较小的核（例如 $3\times 3$ 的权重矩阵），那么当它在输入图像 $f$ 上滑动时，$x-i$ 和 $y-j$ 意味着核的索引与滑窗坐标是反向的。换一种直观的说法：要正确实现数学上的卷积，必须事先将卷积核\textbf{旋转180°}（即颠倒行和列的顺序），再与图像上的局部窗口逐元素相乘并求和。这在信号处理时有明确的物理意义（例如“系统响应的时间反转”），但对于从数据中学习核权重的神经网络而言，这个翻转步骤其实是多余的，因为学到的东西并无区别。

所以为了简便起见，在深度学习的实际操作中，我们采用的运算是下面的\textbf{互相关}运算：
\begin{equation}
    (f \star g)(x, y) = \sum_{i} \sum_{j} f(i, j) \, g(x+i, y+j) \label{eq:cross_correlation}
\end{equation}
互相关与卷积的唯一区别，就在于核 $g$ 的索引是否取负号。在互相关中，核直接以原本的方向滑过图像，窗口内的元素直接相乘求和，完全不需要额外的翻转。

那么，为什么实际使用的不是“纯正的卷积”呢？答案在于：\textbf{核的参数是学出来的}。在训练过程中，无论你事先是否将核翻转，网络都可以通过调整权重来达到完全等价的表达——无非是学习到一个翻转后的核而已。既然反转本身并不带来新的表达能力，却要额外交给计算机一项“翻转坐标”的任务，自然就得不偿失了。省去这步操作不仅降低了计算开销，也让底层实现更加简洁。因此，目前所有主流深度学习框架（例如 PyTorch、TensorFlow）中所谓的“卷积层”（\texttt{Conv2d}），其内部真正执行的正是式 \cref{eq:cross_correlation} 所示的互相关运算\cite{dumoulin2018guideconvolutionarithmeticdeep}。不过，由于历史习惯，整个社区仍然沿用“卷积”这个称呼，后文我们也会继续称其为“卷积层”，而不再每次强调其互相关本质。

\begin{mdframed}[frametitle=例子：二维卷积与互相关]
假设输入是一个 $3\times 3$ 的单通道图像
\[
f = \begin{bmatrix}
a & b & c \\
d & e & f \\
g & h & i
\end{bmatrix}, \quad
\text{核}~~K = \begin{bmatrix}
w_1 & w_2 \\
w_3 & w_4
\end{bmatrix}.
\]
在位置 $(1,1)$ 处，数学上的二维卷积（步长1，不补零）给出：
\[
(f * K)(1,1) = f(0,0)K(1-0,1-0) + f(0,1)K(1-0,1-1) + \cdots 
= a w_4 + b w_3 + d w_2 + e w_1,
\]
这相当于先把 $K$ 旋转180°变成 $ \begin{bmatrix} w_4 & w_3 \\ w_2 & w_1 \end{bmatrix} $，再与对应窗口点乘。
而互相关则直接使用原方向的核：
\[
(f \star K)(1,1) = a w_1 + b w_2 + d w_3 + e w_4.
\]
如果 $w_1,\dots,w_4$ 都是需要优化的自由参数，那么学习到 $\{w_1,\dots,w_4\}$ 与学习到 $\{w_4,w_3,w_2,w_1\}$ 在表达能力上毫无差别。因此用互相关足以替代卷积。
\end{mdframed}
不管叫卷积还是互相关，这种“滑窗加权求和”的结构完美地满足了我们开头提到的两个诉求：
\begin{description}
    \item[平移等变性] 同一个核在图像的不同位置使用\textbf{完全相同的参数}进行运算。这意味着如果猫耳朵在图像中向左平移了一段距离，输出的特征图上相应的激活也会向左平移同样的距离。这正是图\ref{fig:conv-meow}所展示的思想。
    \item[空间局部性] 每一个输出像素只依赖于输入图像中一个很小的邻域（窗口），而不会被远处的像素直接干扰。这使得网络可以专注于局部纹理、边缘等低级特征，再通过多层堆叠逐步组合成全局语义。
\end{description}
正是因为有了这两条优良的性质，卷积（互相关）层才能用微不足道的参数量（一个 $3\times3$ 的核仅有9个参数）替代线性层那种“全连接”所需要的海量权重，从而让 ImageNet-10 这样的大型图像分类问题变得可行。

\subsection{卷积层的前向传播过程}

\begin{figure}
    \centering
    \includestandalone{images/simple-cnn-computation-graph}
    \caption{卷积层的前向传播示意图。输入为一个一阶张量，输出为一个一阶张量，卷积核也是一个一阶张量，卷积核大小为4，步长为1。每个输出元素都是输入张量中连续4个元素与卷积核权重的加权和。}
    \label{fig:conv_forward}
\end{figure}

\paragraph{一层卷积}

在明确了“卷积层实际执行的是互相关运算”之后，我们来形式化地叙述卷积层的前向传播过程：总体而言，一个卷积层一般接受一个三阶张量作为输入，通过一组可学习的卷积核（滤波器）对其进行互相关运算，再加上偏置，最终产生一个新的三阶张量作为输出。

设输入张量为 $\mathbf{X} \in \mathbb{R}^{C_{\text{in}} \times H_{\text{in}} \times W_{\text{in}}}$，其中 $C_{\text{in}}$ 是输入通道数，$H_{\text{in}}$ 和 $W_{\text{in}}$ 分别是输入特征图的高度和宽度。卷积层包含 $C_{\text{out}}$ 个卷积核，每个核的尺寸为 $C_{\text{in}} \times k_h \times k_w$，记作 $\mathbf{W}^{(c)} \in \mathbb{R}^{C_{\text{in}} \times k_h \times k_w}$，$c = 1,\dots,C_{\text{out}}$，这些核各自独立学习不同的权重参数，都要和原张量的所有输入通道进行互相关运算。此外，每个输出通道拥有一个可学习的偏置标量 $b^{(c)}$。

最终，前向传播后得到的输出张量 $\mathbf{Y} \in \mathbb{R}^{C_{\text{out}} \times H_{\text{out}} \times W_{\text{out}}}$ 的第 $c$ 个通道上的位于 $(i,j)$ 的元素为：
\begin{equation}\label{eq:conv_forward}
    \mathbf{Y}_{c,i,j} = \sum_{c'=0}^{C_{\text{in}}-1} \sum_{u=0}^{k_h-1} \sum_{v=0}^{k_w-1} \mathbf{W}^{(c)}_{c',u,v} \, \mathbf{X}_{c', i+u, j+v} \;+\; b^{(c)}.
\end{equation}
这一公式就是带有偏置的互相关运算，实例可见前文的例子。它清楚地表明了每个输出像素是如何通过卷积核在输入特征图上滑动窗口，计算局部加权和来生成的。

一般情况下，卷积核的大小通常取$3\times3$、$5\times5$或$7\times7$，以便在局部邻域内捕捉丰富的空间特征。每个输出通道对应一个卷积核，它们共享相同的输入通道，但各自独立学习不同的权重和偏置，从而能够提取多样化的特征模式。

当 $k_h=k_w=1$ 时，该运算退化为逐点卷积，仅混合通道信息而不涉及空间邻域，但依然遵循同样的形式。

在实际应用中，我们通常还会在空间维度上引入两个整数超参数：\textbf{步长}（stride）和\textbf{零填充}（zero-padding）。步长 $s$ 控制核每次滑动的间隔，零填充 $p$ 则在输入特征图的上下左右各补充 $p$ 行（或列）的零值。引入它们之后，输出特征图的空间尺寸由以下公式决定：
\begin{equation}
    H_{\text{out}} = \left\lfloor\frac{H_{\text{in}} + 2p - k_h}{s}\right\rfloor + 1,\qquad
    W_{\text{out}} = \left\lfloor\frac{W_{\text{in}} + 2p - k_w}{s}\right\rfloor + 1.
\end{equation}
当 $s=1$、$p=0$ 时，输出尺寸恰好为 $(H_{\text{in}}-k_h+1) \times (W_{\text{in}}-k_w+1)$；若希望输出尺寸与输入尺寸保持一致（这在许多现代网络架构中很常见），只需选择 $p = \lfloor k_h/2\rfloor$（假设核尺寸为奇数），同时保持 $s=1$。比如，一个 $3\times3$ 的卷积核配合 $p=1$ 即可得到同样高宽的输出特征图。

\begin{mdframed}[frametitle=例子：单通道图像的卷积层前向传播]

取一幅单通道 $5\times 5$ 的输入图像。使用一个 $3\times3$ 的卷积核 $\mathbf{W}$，偏置 $b=0$，步长 $s=2$，不填充（$p=0$）。则输出尺寸为
\[
H_{\text{out}} = \left\lfloor\frac{5-3}{2}\right\rfloor + 1 = 2,\quad W_{\text{out}}=2.
\]
核从位置 $(0,0)$ 开始，以步长 $2$ 滑动，共计算 $2\times2$ 个输出值。每个输出值都是 $3\times3$ 窗口内图像块与 $\mathbf{W}$ 的内积。
\end{mdframed}

细心的读者或许会问：前文强调卷积层具有\textbf{平移等变性}，这在前向传播公式中如何体现？事实上，对于同一个核，无论窗口滑动到图像的哪一个位置，卷积核的参数 $\mathbf{W}^{(c)}$ 和偏置 $b^{(c)}$ 都是\textbf{完全相同}的。换句话说，公式 \cref{eq:conv_forward} 的权重与绝对坐标 $(i,j)$ 无关，只依赖于相对位置 $(u,v)$。因此，输入信号的平移将等比地映射到输出信号的平移，这正是平移等变性的数学表达。而空间局部性则体现在求和仅限于一个 $k_h \times k_w$ 的小窗口内，远超窗口范围的像素在前向传播中完全不起作用。

\paragraph{多层卷积} 多层卷积堆叠时，前向传播就是逐层执行上述运算。对于某一层输出特征图上的某一个像素（可以看作是一个神经元），我们用“感受野”来称呼它在原始输入图像上对应的区域大小，或者说“和多大块的像素有关”。随着卷积层的堆叠，感受野会逐渐扩大：
\[ \text{感受野} = \text{上一层感受野} + (k_h - 1) \times \text{步长} \]
这意味着，网络的高层神经元能够“看到”输入图像中更大范围的像素信息，从而捕捉到更抽象、更全局的特征。

就这样逐层学习，第一层直接处理原始像素，学习边、角、颜色等低级特征；第二层在第一层输出的特征图上继续滑窗，组合出纹理、形状等更复杂的模式；随着层数加深，感受野逐步扩大，网络最终能够在高层特征图上捕捉到足以支撑分类的全局语义信息。这种逐层抽象的机制，正是卷积神经网络强大表达能力的核心来源。

但必须说明的是，单纯的卷积层不能做到给张量降阶：那是线性层才能做的事情。所以，我们一般在卷积层之后添加一个全连接层，将卷积层输出的三阶张量展平成一阶张量，再输入到全连接层中进行最终的分类或回归任务。

\paragraph{卷积层的参数量} 既然构建出了卷积层，那么我们自然好奇它有多少参数可以学习。

卷积层的参数量完全决定于卷积核的数量和大小。因为每一个卷积核都是一个三阶张量，还知道这些张量的数量、每个张量的大小，以及每个输出通道都有一个偏置参数，于是就可以轻易的求出卷积层有多少参数了。

设输入通道数为 $C_{\text{in}}$，输出通道数为 $C_{\text{out}}$，卷积核的尺寸为 $k_h \times k_w$，则卷积层的总参数量为：
\[\text{参数量} = C_{\text{out}} \times (C_{\text{in}} \times k_h \times k_w + 1),\]
其中 $+1$ 是每个输出通道对应的偏置参数。相比于全连接层，卷积层的参数量通常要小得多，因为它们共享权重并且只在局部区域内进行计算，这使得卷积神经网络在处理高维图像数据时更加高效和可扩展。

\paragraph{特殊的卷积：$1\times 1$卷积} 当卷积核的空间尺寸为 $1\times 1$ 时，卷积层退化为逐点卷积。此时，每个输出像素仅依赖于输入像素在同一位置的所有通道信息，而不涉及空间邻域。这种操作常用于通道数的压缩或扩展，例如在网络中间层进行特征图的降维或升维。尽管没有空间感受野，但 $1\times 1$ 卷积仍然保留了平移等变性和参数共享的优势。

\subsection{卷积层的反向传播}

刚刚我们研究了卷积层的前向传播，现在来研究一下卷积层的反向传播。由于链式法则，我们只需要研究卷积层的输出对输入的梯度。

对于卷积层而言，好消息是它的局部操作极其规整——仍是一系列互相关（或与之等价的卷积）形式，这使得我们可以用与前向传播几乎相同的高效实现来完成反向传播。

为便于表述，我们暂时略去填充和步长，聚焦于最朴素的情形：输入为 $\mathbf{X}\in\mathbb{R}^{C_{\text{in}}\times H_{\text{in}}\times W_{\text{in}}}$，核为 $\mathbf{W}^{(c)}\in\mathbb{R}^{C_{\text{in}}\times k_h\times k_w}$（$c=0,\dots,C_{\text{out}}-1$），输出为 $\mathbf{Y}$，其第 $c$ 个通道上位置 $(i,j)$ 的元素由式 \cref{eq:conv_forward} 给出：
\[
\mathbf{Y}_{c,i,j} = \sum_{c'}\sum_{u}\sum_{v} \mathbf{W}^{(c)}_{c',u,v} \, \mathbf{X}_{c', i+u, j+v} + b^{(c)}.
\]
记总损失为 $\mathcal{L}$，我们已经得到上游梯度 $\frac{\partial\mathcal{L}}{\partial \mathbf{Y}}$。现在需要求 $\frac{\partial\mathcal{L}}{\partial \mathbf{X}}$，$\frac{\partial\mathcal{L}}{\partial \mathbf{W}^{(c)}}$ 和 $\frac{\partial\mathcal{L}}{\partial b^{(c)}}$。

\paragraph{对偏置的梯度}
偏置 $b^{(c)}$ 在输出通道 $c$ 的每一个空间位置上都被加了同样的值，因此
\[
\frac{\partial\mathcal{L}}{\partial b^{(c)}} = \sum_{i=0}^{H_{\text{out}}-1}\sum_{j=0}^{W_{\text{out}}-1} \frac{\partial\mathcal{L}}{\partial \mathbf{Y}_{c,i,j}}.
\]
简言之，就是把输出梯度对该通道的所有位置求和。

\paragraph{对核权重的梯度}
考虑权重 $\mathbf{W}^{(c)}_{c',u,v}$，它在计算 $\mathbf{Y}_{c,i,j}$ 时与输入 $\mathbf{X}_{c', i+u, j+v}$ 相乘。对于固定的 $u,v$，该权重被所有合法的输出位置 $(i,j)$ 共享。因此，由多元链式法则：
\[
\frac{\partial\mathcal{L}}{\partial \mathbf{W}^{(c)}_{c',u,v}} = \sum_{i=0}^{H_{\text{out}}-1}\sum_{j=0}^{W_{\text{out}}-1} \frac{\partial\mathcal{L}}{\partial \mathbf{Y}_{c,i,j}} \cdot \frac{\partial \mathbf{Y}_{c,i,j}}{\partial \mathbf{W}^{(c)}_{c',u,v}}
= \sum_{i,j} \frac{\partial\mathcal{L}}{\partial \mathbf{Y}_{c,i,j}} \; \mathbf{X}_{c', i+u, j+v}.
\]
观察上式，它正是将输出梯度 $\frac{\partial\mathcal{L}}{\partial \mathbf{Y}_{c,:,:}}$ 作为“核”，在输入 $\mathbf{X}_{c',:,:}$ 上执行一次互相关运算——与式 \cref{eq:cross_correlation} 完全相同，只不过现在“图像”是输入，“核”是输出梯度。换言之，\textbf{权重的梯度可以通过前向传播中一模一样的互相关操作得到}。这让实际编程实现异常简洁：复用同一套滑窗乘加代码即可。

\paragraph{对输入的梯度}
最后，也是最需要留意的，是损失对输入 $\mathbf{X}$ 的梯度。一个输入像素 $\mathbf{X}_{c',p,q}$ 会被所有满足 $p = i+u$、$q = j+v$ 的输出位置所使用。每个这样的输出位置贡献一条梯度路径，携带上游梯度的相应分量和对应的核权重。因此：
\[
\frac{\partial\mathcal{L}}{\partial \mathbf{X}_{c',p,q}} = \sum_{c=0}^{C_{\text{out}}-1} \sum_{u=0}^{k_h-1}\sum_{v=0}^{k_w-1} 
\frac{\partial\mathcal{L}}{\partial \mathbf{Y}_{c,\,p-u,\;q-v}} \; \mathbf{W}^{(c)}_{c',u,v},
\]
其中约定当索引 $p-u$ 或 $q-v$ 超出有效输出范围时，对应的上游梯度取 0（即不考虑填充效应下的边界）。仔细观察求和形式，内层对于 $u,v$ 的累加可以看作一个与翻转后的核的运算：若定义核的翻转 $\widetilde{\mathbf{W}}^{(c)}_{c',u,v} = \mathbf{W}^{(c)}_{c',\,k_h-1-u,\,k_w-1-v}$，则上式等价于
\[
\frac{\partial\mathcal{L}}{\partial \mathbf{X}_{c',p,q}} = \sum_{c=0}^{C_{\text{out}}-1} \sum_{u=0}^{k_h-1}\sum_{v=0}^{k_w-1} 
\frac{\partial\mathcal{L}}{\partial \mathbf{Y}_{c,\,p-u,\;q-v}} \; \widetilde{\mathbf{W}}^{(c)}_{c',\,k_h-1-u,\,k_w-1-v}.
\]
进一步可将指标平移，使其变成标准的互相关形式，但更自然的观点是：\textbf{输入梯度等于输出梯度与翻转后的核进行的数学卷积}。这正是为什么在教科书和早期代码中，人们常说“反向传播是卷积操作”——因为它确实严格对应于数学上的卷积（核被翻转）。而在现代框架统一使用互相关的语境下，只需记住：求输入梯度的操作，本质上是用核的\textbf{转置}（空间方向上下左右颠倒）与输出梯度做互相关，或者在补充零填充之后执行一个“全卷积”。

\begin{mdframed}[frametitle=例子：单通道小图的卷积反向传播（不填充，步长1）]

考虑输入为 $3\times3$ 的单通道图像：
\[
\mathbf{X} = \begin{bmatrix}
x_{0,0} & x_{0,1} & x_{0,2} \\
x_{1,0} & x_{1,1} & x_{1,2} \\
x_{2,0} & x_{2,1} & x_{2,2}
\end{bmatrix}, \qquad
\mathbf{W} = \begin{bmatrix}
w_{0,0} & w_{0,1} \\
w_{1,0} & w_{1,1}
\end{bmatrix}.
\]
偏置为0，步长1，输出 $\mathbf{Y}$ 为 $2\times2$ 矩阵，其中
\[
Y_{i,j} = w_{0,0}x_{i,j} + w_{0,1}x_{i,j+1} + w_{1,0}x_{i+1,j} + w_{1,1}x_{i+1,j+1}.
\]
设上游梯度 $\frac{\partial\mathcal{L}}{\partial \mathbf{Y}} = \begin{bmatrix} \delta_{0,0} & \delta_{0,1} \\ \delta_{1,0} & \delta_{1,1} \end{bmatrix}$。
则权重的梯度为
\[
\frac{\partial\mathcal{L}}{\partial w_{u,v}} = \sum_{i=0}^1\sum_{j=0}^1 \delta_{i,j} \cdot \mathbf{X}_{i+u, j+v}.
\]
对于输入 $\mathbf{X}$，例如 $x_{1,1}$（中心点）出现在 $Y_{0,0}$（系数 $w_{1,1}$）、$Y_{0,1}$（系数 $w_{1,0}$）、$Y_{1,0}$（系数 $w_{0,1}$）、$Y_{1,1}$（系数 $w_{0,0}$）中，因此
\[
\frac{\partial\mathcal{L}}{\partial x_{1,1}} = \delta_{0,0}w_{1,1} + \delta_{0,1}w_{1,0} + \delta_{1,0}w_{0,1} + \delta_{1,1}w_{0,0},
\]
恰好对应将 $\mathbf{W}$ 上下左右翻转后（$\widetilde{\mathbf{W}} = \begin{bmatrix}w_{1,1}&w_{1,0}\\ w_{0,1}&w_{0,0}\end{bmatrix}$）与 $\delta$ 做互相关（或说 $\delta$ 与 $\widetilde{\mathbf{W}}$ 的数学卷积）得到的中心位置的值。
\end{mdframed}

现实实现中，处理填充和步长会让索引边界稍显繁琐，但核心思想不变：\textbf{卷积层的反向传播依旧是互相关/卷积运算}。为追求并行效率，常用方式是将局部窗口展开成矩阵，此时前向传播变为矩阵乘法，反向传播则自动化为矩阵乘法的转置——梯度对输入的传播将用到核矩阵的转置乘以上游梯度，而对权重的梯度则是输入矩阵的转置乘以上游梯度。这正是我们熟悉的“全连接层的反向传播”，只不过矩阵的形状由卷积参数（核大小、步长、填充）决定。

\subsection{特殊的卷积：滤波器}

当卷积核固定时，卷积层就不再是可学习的，而是确定性的。通常把这样的卷积核叫做“滤波器”，它们在图像处理中有着悠久的历史。滤波器的作用是对图像进行特定的处理，例如边缘检测、模糊处理、锐化等。常见的滤波器包括Sobel算子、Laplacian算子、高斯模糊等。

在深度神经网络兴起之前，计算机视觉几乎完全依赖于手工设计的滤波器来提取特征。这种固定核的思想其实与我们之前讨论的卷积层一脉相承：使用一个小型窗口在整个图像上滑动，通过对局部像素的加权求和实现某种变换。唯一的区别在于，这里的权重由人类根据数学或物理原理预先设定，而非从数据中通过梯度下降学习得到。本节我们简要介绍几种经典滤波器，并讨论它们与可学习卷积层的关系。

\paragraph{平滑与高斯模糊}
图像中不可避免会存在噪声，平滑滤波器通过将每个像素替换为其邻域内像素的加权平均值来抑制噪声。最简单的平滑核是“盒式滤波器”：
\[
K_{\text{box}} = \frac{1}{9}\begin{bmatrix}
1 & 1 & 1 \\
1 & 1 & 1 \\
1 & 1 & 1
\end{bmatrix}.
\]
它赋予邻域内每个像素相同的权重。更自然的平滑是基于高斯函数生成的核，离中心越近的像素权重越大，符合“近者影响大”的直觉。例如，$3\times3$ 高斯核（$\sigma\approx 0.85$）通常近似为：
\[
K_{\text{Gauss}} \approx \frac{1}{16}\begin{bmatrix}
1 & 2 & 1 \\
2 & 4 & 2 \\
1 & 2 & 1
\end{bmatrix}.
\]
用这样的核对图像做互相关，会抹去细节、突出大面积均匀区域，也就是“模糊”了图像。

\paragraph{边缘检测：Sobel算子}
边缘是图像中亮度发生剧烈变化的位置，对应于一阶导数的极值或二阶导数的过零点。Sobel算子通过计算水平和垂直两个方向的离散导数来检测边缘。其水平边缘检测核（计算垂直梯度）和垂直边缘检测核（计算水平梯度）分别为：
\[
G_x = \begin{bmatrix}
-1 & 0 & 1 \\
-2 & 0 & 2 \\
-1 & 0 & 1
\end{bmatrix},\qquad
G_y = \begin{bmatrix}
-1 & -2 & -1 \\
 0 &  0 &  0 \\
 1 &  2 &  1
\end{bmatrix}.
\]
对于图像 $I$，分别用 $G_x$ 和 $G_y$ 与之做互相关，得到水平梯度图 $I_x$ 和垂直梯度图 $I_y$。梯度的幅值 $M = \sqrt{I_x^2 + I_y^2}$ 可以近似反映边缘强度，方向 $\theta = \arctan(I_y / I_x)$ 给出了边缘的法向。Sobel算子兼具平滑（中间行或列权重较大）和差分作用，对噪声有一定抑制能力，是图像处理中最常用的边缘检测器之一。

\paragraph{锐化与拉普拉斯算子}
锐化的目的是增强边缘和细节，本质上是对图像的平滑版本做减法（非锐化掩模）或直接用拉普拉斯算子提取二阶导数信息。拉普拉斯算子是最简单的各向同性二阶微分算子，其常用离散核为：
\[
\nabla^2 \approx \begin{bmatrix}
0 &  1 & 0 \\
1 & -4 & 1 \\
0 &  1 & 0
\end{bmatrix}\quad\text{或}\quad
\begin{bmatrix}
1 &  1 & 1 \\
1 & -8 & 1 \\
1 &  1 & 1
\end{bmatrix}.
\]
将拉普拉斯响应叠加回原图像（即 $I_{\text{sharp}} = I - \alpha \nabla^2 I$，$\alpha$ 为一个小正数），即可实现图像锐化。

\begin{figure}
    \centering
    \includegraphics[width=\textwidth]{conv-meow-handled.png}
    \caption{使用固定滤波器处理图像}
\end{figure}

\paragraph{从固定滤波器到可学习卷积}
在深度学习的语境下，我们不再手工设计这些核，而是让网络自动学习最适合当前任务的滤波器。然而，理解经典滤波器仍然很有价值：
\begin{enumerate}
    \item 它们常常作为\textbf{初始化}的参考。例如，某些工作使用Gabor滤波器初始化第一层卷积核，以模仿生物视觉皮层的响应特性\cite{zeiler2013visualizingunderstandingconvolutionalnetworks}。
    \item 当我们可视化训练好的卷积层权重时，常常会发现低层的核自动演化成了类似边缘检测器、颜色斑块探测器等模式，与Sobel、Gabor等滤波器非常相似。这说明网络自发地发现了“边缘检测”这一有效的底层特征表达。
    \item 在一些应用中（如图像预处理、风格迁移的感知损失计算），我们仍然故意使用固定且不可训练的滤波器来提取特定信息，而不会让它们参与梯度更新。只不过在当今端到端的深度学习框架中，这种做法越来越少见。
\end{enumerate}

\subsection{\mintinline{python}{PyTorch}中的卷积层}
在 \mintinline{python}{PyTorch} 中，卷积层由 \mintinline{python}{torch.nn.Conv2d} 类实现。其构造函数的主要参数包括：
\begin{itemize}
    \item \mintinline{python}{in_channels}：输入特征图的通道数 $C_{\text{in}}$。
    \item \mintinline{python}{out_channels}：输出特征图的通道数 $C_{\text{out}}$。
    \item \mintinline{python}{kernel_size}：卷积核的空间尺寸 $(k_h, k_w)$，可以是单个整数（表示方形核）或一个二元组。
    \item \mintinline{python}{stride}：步长 $s$，默认为1。
    \item \mintinline{python}{padding}：零填充 $p$，默认为0。
    \item \mintinline{python}{bias}：是否使用偏置，默认为True。
\end{itemize}
使用示例：
\begin{minted}{python}
import torch
import torch.nn as nn

conv_layer = nn.Conv2d(in_channels=3, out_channels=16, kernel_size=3, stride=1, padding=1, bias=True)
input_tensor = torch.randn(1, 3, 32, 32)  # Batch size 1, 3 channels, 32x32 image
output_tensor = conv_layer(input_tensor)
print(output_tensor.shape)  # Should be (1, 16, 32, 32) due to padding=1
\end{minted}

这样的卷积层可以放在神经网络中，例如：
\begin{minted}{python}
class SimpleCNN(nn.Module):
    def __init__(self):
        super(SimpleCNN, self).__init__()
        self.conv1 = nn.Conv2d(3, 16, 3, stride=1, padding=1)
        self.conv2 = nn.Conv2d(16, 32, 3, stride=1, padding=1)
        self.fc = nn.Linear(32 * 32 * 32, 10)  # Assuming input images are 32x32
    def forward(self, x):
        x = torch.relu(self.conv1(x))
        x = torch.relu(self.conv2(x))
        x = x.view(x.size(0), -1)  # Flatten
        x = self.fc(x)
        return x
\end{minted}

\section{池化层}

卷积层虽然大大降低了参数量，但仍然需要一个核函数来进行局部加权映射，这个核函数是可学习的。既然可以学习，那就需要算力。对于多数任务而言，我们并不需要那么精确的局部映射，而是希望在空间上进行降采样，保留主要特征，此时卷积层就“太重”了。而池化层，正是为此而设计的。

\subsection{池化层及其前向传播}

池化层的思路和卷积层差不多：在输入特征图上滑动一个窗口，计算窗口内的某种统计量（如最大值、平均值）作为输出。不同的是，池化层没有可学习的参数，它只执行固定的操作，因此计算开销更小。

最常见的池化方式就是求最大值和平均值，对应叫做“最大池化”和“平均池化”。设输入特征图为 $\mathbf{X} \in \mathbb{R}^{C \times H \times W}$，池化窗口大小为 $k_h \times k_w$，步长为 $s$，则输出特征图 $\mathbf{Y} \in \mathbb{R}^{C \times H_{\text{out}} \times W_{\text{out}}}$ 的每个元素计算如下：
\[
\mathbf{Y}_{c,i,j} = \text{pool}\left(\mathbf{X}_{c, i \cdot s : i \cdot s + k_h, j \cdot s : j \cdot s + k_w}\right),
\]
其中 $\text{pool}$ 可以是 $\max$ 或 $\text{mean}$，表示在窗口内取最大值或平均值。可以看到这个公式和卷积层如出一辙。

所以自然可以引入池化窗口的填充和步长这两个超参数，最终输出特征图的空间尺寸为：
\[
H_{\text{out}} = \left\lfloor\frac{H + 2p - k_h}{s}\right\rfloor + 1,\qquad
W_{\text{out}} = \left\lfloor\frac{W + 2p - k_w}{s}\right\rfloor + 1.
\]

池化层没有可以学习的参数，因此它不会像卷积层那样通过梯度下降来更新权重。它的主要作用是降低特征图的空间尺寸，从而减少计算量和内存占用，同时也能在一定程度上提供平移等变性和空间局部性。

一般不会堆叠多个池化层。我们更多的是将池化层放在卷积层的后面，将较繁琐的卷积层输出的特征图进行降采样，减少后续层的计算负担。

\subsection{池化层的反向传播}

池化层虽然不参与学习，但依然是神经网络的组成部分（类似激活函数），所以依然需要反向传播、计算梯度。

对于最大池化，反向传播的关键是找到在前向传播中被选为最大值的那个输入位置。设上游梯度为 $\frac{\partial \mathcal{L}}{\partial \mathbf{Y}}$，则对于每个池化窗口，只有最大值对应的输入位置会接收到梯度，其余位置的梯度为零。

对于平均池化，反向传播则是将上游梯度均匀分配到池化窗口内的所有输入位置。设池化窗口大小为 $k_h \times k_w$，则每个输入位置接收到的梯度为：
\[
\frac{\partial \mathcal{L}}{\partial \mathbf{X}_{c, i, j}} = \frac{1}{k_h \cdot k_w} \cdot \frac{\partial \mathcal{L}}{\partial \mathbf{Y}_{c, \lfloor i/s \rfloor, \lfloor j/s \rfloor}}.
\]


\subsection{\mintinline{python}{PyTorch}中的池化层}

在 \mintinline{python}{PyTorch} 中，池化层由 \mintinline{python}{torch.nn.MaxPool2d} 和 \mintinline{python}{torch.nn.AvgPool2d} 类实现。其构造函数的主要参数包括：

\begin{itemize}
    \item \mintinline{python}{kernel_size}：池化窗口的空间尺寸 $(k_h, k_w)$，可以是单个整数（表示方形窗口）或一个二元组。
    \item \mintinline{python}{stride}：步长 $s$，默认为池化窗口大小。
    \item \mintinline{python}{padding}：零填充 $p$，默认为0。
\end{itemize}
使用示例：
\begin{minted}{python}
import torch
import torch.nn as nn
max_pool = nn.MaxPool2d(kernel_size=2, stride=2)
input_tensor = torch.randn(1, 3, 32, 32)  # Batch size 1, 3 channels, 32x32 image
output_tensor = max_pool(input_tensor)
print(output_tensor.shape)  # Should be (1, 3, 16, 16) due to stride=2
\end{minted}

在神经网络中一般放在卷积层的后面，作为特征图的降采样操作。例如：
\begin{minted}{python}
class CNNWithPooling(nn.Module):
    def __init__(self):
        super(CNNWithPooling, self).__init__()
        self.conv1 = nn.Conv2d(3, 16, 3, stride=1, padding=1)
        self.pool = nn.MaxPool2d(2, stride=2)
        self.conv2 = nn.Conv2d(16, 32, 3, stride=1, padding=1)
        self.fc = nn.Linear(32 * 8 * 8, 10)  # Assuming input images are 32x32
    def forward(self, x):
        x = torch.relu(self.conv1(x))
        x = self.pool(x)
        x = torch.relu(self.conv2(x))
        x = self.pool(x)
        x = x.view(x.size(0), -1)  # Flatten
        x = self.fc(x)
        return x
\end{minted}

\section{上采样层}

卷积层和池化层有一个共同的特点：倾向于将输出特征图的空间尺寸缩小，从而提取出更抽象的特征。然而，在某些任务中（如图像生成、语义分割），我们需要将特征图的空间尺寸恢复到原始大小或更大，这就需要上采样（和缩小尺寸的“下采样”对应）。

常见的上采样方法包括最近邻插值、线性插值和转置卷积。

\subsection{最近邻插值与线性插值}

最近邻插值可以看作是受到了平均池化的启发：它将每个像素复制到其邻域内，从而扩大图像尺寸。具体来说，若我们希望将图像的宽高各扩大 $s$ 倍，则每个像素 $(i,j)$ 的值会被复制到输出图像中 $(si, sj)$ 到 $(si+s-1, sj+s-1)$ 的区域内。这是非常朴素的做法。

线性插值则更为平滑，它在插值过程中考虑了邻近像素的加权平均，从而生成过渡更自然的图像。对于二维图像，常用的是双线性插值，它在水平方向和垂直方向上分别进行线性插值。线性插值的公式则是著名的$x = (1-t)x_0 + t x_1$，其中$t$是插值位置的相对距离。

这两种插值手段都是不依赖于学习的，它们只是根据已有像素值进行计算，因此不涉及参数更新。

\subsection{转置卷积}

最近邻插值和线性插值虽然简单高效，但它们完全不涉及可学习的参数，因此无法根据任务自适应地调整上采样策略。在深度学习的许多任务中（如语义分割、图像生成），我们希望上采样本身也是网络学习的一部分。这就引出了\textbf{转置卷积}（transposed convolution），它更广为人知但容易引起误解的名字是“反卷积”（deconvolution）。

\paragraph{为什么不该叫“反卷积”}

“反卷积”这个词汇原本在信号处理中有严格的定义：给定一个线性系统的输出和系统响应，求解原始输入信号的过程。它是卷积的\textbf{逆运算}，在数学上是唯一确定的。然而，我们在深度学习中使用的“反卷积层”根本不是在求卷积的逆——它本身也是一次卷积（互相关）操作，只不过步长和填充的配置使得输出尺寸大于输入尺寸。更准确地说，它在反向传播中计算梯度的方式恰好是正向卷积的转置，因此文献中逐渐以\textbf{转置卷积}作为正式名称\cite{dumoulin2018guideconvolutionarithmeticdeep}。为免混淆，后文直接使用这一术语。

\paragraph{前向传播：用分数步长实现尺寸放大}

回想一下，普通卷积（步长 $s\ge 1$，不填充）会将输入特征图缩小大约 $s$ 倍。转置卷积的做法颇为巧妙：它\textbf{先在输入像素之间插入零值}（使特征图“膨胀”），再以步长 $1$ 对该膨胀后的图像执行普通卷积。这样一来，输出特征图的尺寸不但不会缩小，反而会被放大。

具体而言，设输入为 $\mathbf{X}\in\mathbb{R}^{H_{\text{in}}\times W_{\text{in}}}$（为简化符号，暂考虑单通道），核为 $\mathbf{W}\in\mathbb{R}^{k_h\times k_w}$，步长为 $s$（这里指上采样倍率），零填充为 $p$。转置卷积的前向传播可以拆解为两步：
\begin{enumerate}
    \item \textbf{零插入}：在 $\mathbf{X}$ 的每两个相邻行之间插入 $s-1$ 行全零，每两个相邻列之间插入 $s-1$ 列全零。记膨胀后的特征图为 $\tilde{\mathbf{X}}$，其尺寸为
    \[
    \tilde{H} = s(H_{\text{in}}-1) + 1,\quad \tilde{W} = s(W_{\text{in}}-1) + 1.
    \]
    \item \textbf{普通卷积}：以步长 $1$、零填充 $p$，用核 $\mathbf{W}$ 对 $\tilde{\mathbf{X}}$ 执行互相关运算，得到输出 $\mathbf{Y}$。
\end{enumerate}
将这两步合并，输出尺寸的公式为：
\begin{equation}
    H_{\text{out}} = s(H_{\text{in}}-1) - 2p + k_h,\qquad
    W_{\text{out}} = s(W_{\text{in}}-1) - 2p + k_w. \label{eq:transposed_conv_size}
\end{equation}
当 $s>1$ 时，输出尺寸明显大于输入尺寸，这正是我们需要的上采样效果。

\begin{mdframed}[frametitle=例子：最简单的转置卷积]

输入为 $2\times2$ 的单通道矩阵
\[
\mathbf{X} = \begin{bmatrix}
a & b \\
c & d
\end{bmatrix},
\]
核为 $3\times3$，步长 $s=2$，不填充（$p=0$）。按照公式 \eqref{eq:transposed_conv_size}，输出尺寸为
\[
H_{\text{out}} = 2\times(2-1) - 0 + 3 = 5,\quad W_{\text{out}} = 5.
\]
零插入后，$\tilde{\mathbf{X}}$ 变为 $3\times3$，其中仅有原输入位置的四个值为非零，其余皆为 $0$。随后用 $3\times3$ 核对这个稀疏矩阵以步长 $1$ 滑动做互相关，每个输出位置将只“看到”落入核窗口内的几个非零输入值及其对应的权重，最终生成 $5\times5$ 的输出。这个过程可以直观地理解为：\textbf{每一个输入标量被核的所有元素加权后“散布”到输出的一个对应区域，相邻输入散布的区域相互重叠叠加，构成了最终的输出。}
\end{mdframed}

\paragraph{反向传播：本质是正向卷积}

转置卷积之所以跟“转置”挂钩，是因为从梯度流的角度看，它与普通卷积形成完美的对称。回顾上一节卷积层反向传播的结论：
\begin{itemize}
    \item 正向普通卷积（步长 $s$，填充 $p$）的反向传播，等价于用翻转后的核对输出梯度做步长 $1$、填充 $k_h-1-p$ 的卷积（或将梯度先膨胀再与核做互相关）。
    \item 反过来，转置卷积（零插入 $+$ 步长1卷积）的反向传播，恰恰就是一个\textbf{步长为 $s$、填充为 $p$ 的普通卷积}，直接用同一核在输出梯度上滑动。
\end{itemize}
因此，如果我们把卷积层的前向操作记为矩阵乘法 $\mathbf{Y} = \mathbf{C}\mathbf{X}$（其中 $\mathbf{C}$ 是由卷积核和步长、填充决定的稀疏矩阵），那么普通卷积的反向传播是左乘 $\mathbf{C}^\top$，而转置卷积的前向传播正是左乘 $\mathbf{C}^\top$，反向传播再左乘 $(\mathbf{C}^\top)^\top = \mathbf{C}$。两者恰好在矩阵运算的意义上互为转置，这就是名称的真正由来\cite{dumoulin2018guideconvolutionarithmeticdeep}。

\paragraph{棋盘效应与缓解手段}

尽管转置卷积被广泛使用，它有一个著名的问题：\textbf{棋盘状伪影}（checkerboard artifacts）。当核的尺寸不能被步长整除时，零插入后的稀疏输入会导致输出特征图中某些像素被覆盖的次数多于其他像素，形成不均匀的重叠。这种不均匀性在视觉上表现为类似棋盘的明暗交替图案，在高分辨率图像生成中尤为明显\cite{2016Deconvolution}。

缓解棋盘效应的方法包括：
\begin{enumerate}
    \item \textbf{选择核大小可被步长整除}：例如步长为 $2$ 时，核大小选择 $4\times4$ 而非 $3\times3$。这样每个输入像素散布的区域恰好可以无缝拼接，避免重叠不均匀。
    \item \textbf{后处理}：在转置卷积之后额外添加一层普通卷积或平滑滤波，以抹平不均匀的幅度。
    \item \textbf{改用“插值+卷积”}：先用最近邻或双线性插值将特征图放大到目标尺寸，再用普通卷积进行特征精炼。这种方法近年来逐渐取代转置卷积成为上采样层的主流实现，因为它完全避免了棋盘效应，且计算量相当。
\end{enumerate}

尽管存在这些注意事项，转置卷积依然深刻体现了“让网络自己学会如何上采样”这一思想。它和普通卷积共享同一套核参数，使得下采样和上采样路径可以在编码器-解码器结构（如U-Net\cite{ronneberger2015unetconvolutionalnetworksbiomedical}）中形成优雅的对应。理解它的运作机制和命名渊源，也能帮助我们更清晰地读懂各类网络结构图中的“ConvTranspose2d”到底在做什么。

\section{卷积神经网络在计算机视觉领域的典型应用}

计算机视觉，指的是让计算机理解和分析图像或视频内容的技术领域。

卷积神经网络是处理图像的利器。自从LeNet-5\cite{1998Gradient}在手写数字识别上取得突破以来，卷积神经网络的架构不断演进，涌现出一系列经典模型，解决了许多困难的视觉任务。CNN 被广泛应用于图像分类、目标检测、语义分割、实例分割、人脸识别、图像生成等任务。其优势在于能够自动学习图像的层次化特征表示，从而减少对手工特征工程的依赖。

\subsection{图像分类}

图像分类是计算机视觉中最基础的任务之一，目标是将输入图像分配到预定义的类别标签中。卷积神经网络在图像分类任务中表现出色，能够自动提取图像的空间特征，并通过多层非线性变换实现复杂的模式识别。

\paragraph{LeNet-5}

LeNet-5~\cite{1998Gradient} 诞生于1998年，是早期最著名的卷积神经网络之一，最初用于银行支票的手写数字识别。它的结构极为简洁：输入为 $32\times32$ 的灰度图像，依次经过两个卷积层（每个卷积层后接平均池化层），再经过两个全连接层，最终输出10个类别。两个卷积核的尺寸均为 $5\times5$，分别输出6个和16个特征图；池化窗口为 $2\times2$，采用平均池化并跟随一个可训练的线性缩放与偏置。全连接层将特征图展平后依次映射到120维和84维，最后通过一个径向基函数（RBF）输出层进行分类。

LeNet-5 第一次系统地演示了“卷积–池化–全连接”这种层次化特征提取范式的有效性，用很少的参数量（约6万）在手写数字任务上达到了与当时最优秀的特征工程方法相当甚至更优的水平。它的成功开启了卷积神经网络的大门，证明局部连接和权值共享是处理图像的有力手段。

不过，受限于当时的计算能力与数据规模，LeNet-5 无法扩展到更大、更复杂的图像问题。它缺乏现代训练技巧（如ReLU激活函数、批量归一化、Dropout），在深层网络上极易陷入梯度消失。此外，它的输出层使用了特定于任务的设计（RBF），不如后来的Softmax灵活通用。

\paragraph{AlexNet}

2012年，AlexNet~\cite{2012ImageNet} 以显著优势夺得 ImageNet 大规模视觉识别挑战赛（ILSVRC）的冠军，正式宣告深度学习时代的到来。AlexNet 的网络结构整体上可看作 LeNet-5 的放大与现代化版本：输入为 $224\times224\times3$ 的彩色图像，包含5个卷积层和3个全连接层，最后一个全连接层输出1000类分布。关键改动包括：首次采用 ReLU 激活函数替代 Sigmoid，极大缓解了梯度消失；引入重叠的最大池化（步长小于窗口尺寸）以增加特征多样性；在全连接层使用 Dropout 抑制过拟合；并使用本地响应归一化（Local Response Normalization）增强侧抑制效应。由于当时的 GPU 显存仍很有限，训练时还将模型拆分到两块 GPU 上并行计算，仅在特定层间交换特征图。

AlexNet 凭借一系列训练技巧将深度卷积网络推向了实用化，首次证明了在百万量级的数据集上训练深层网络不仅可行，而且效果远超传统手工特征+浅层分类器的组合。它的成功直接引发了计算机视觉领域的“深度学习革命”。

但该网络的超参数选择（如核大小、步长等）带有较强的经验色彩，缺乏系统性的设计原则。较大的全连接层（如第一层全连接有4096个神经元）导致参数量膨胀至约6000万，占用了大量显存，且容易过拟合。此外，本地响应归一化后来被证明效果有限，基本被批量归一化所取代。

\paragraph{VGGNet}

VGGNet~\cite{simonyan2015deepconvolutionalnetworkslargescale} 由牛津大学视觉几何组（Visual Geometry Group）于2014年提出，其核心思想是“用更深的网络，但仅使用小尺寸的卷积核”。最常用的变体 VGG-16 包含13个 $3\times3$ 卷积层和3个全连接层，所有卷积层均使用大小为 $3\times3$、步长为1的卷积核，搭配同样的 padding 以保持特征图尺寸不变；池化层固定为 $2\times2$ 的最大池化、步长2，用于下采样。通道数从64开始，每次空间尺寸减半后翻倍，直至512。这种极度规整的结构使网络可以轻松地加深到16或19层。

小卷积核的堆叠带来了两大好处：第一，两层 $3\times3$ 卷积的有效感受野等效于一层 $5\times5$ 卷积，但参数量更少（$2\times 3^2 < 5^2$）且包含更多非线性变换；第二，结构极度规则，易于理解和迁移。VGGNet 在 ImageNet 上取得了当时顶尖的性能，其预训练模型至今仍被广泛用于特征提取和风格迁移等任务。

但全连接层的存在导致 VGGNet 的参数量极为庞大（VGG-16 约1.38亿参数），且绝大部分参数集中在靠后的全连接层，既消耗大量存储与计算资源，又容易过拟合。另外，单纯的顺序堆叠没有解决梯度消失的问题，层数再增加时性能反而退化。

\paragraph{GoogLeNet}

同样在2014年，Google 提出了 GoogLeNet~\cite{szegedy2014goingdeeperconvolutions}（亦称 Inception-v1），它以完全不同的思路解决“如何加深网络”的问题。GoogLeNet 引入了 \textbf{Inception 模块}，其基本思想是“让网络自己决定用多大的卷积核”：在同一层中，输入特征图同时通过 $1\times1$ 卷积、$3\times3$ 卷积、$5\times5$ 卷积和 $3\times3$ 最大池化四条并行的路径，所有输出在通道维度上拼接。为了控制计算量，在较大的卷积核前以及池化路径后均插入了 $1\times1$ 卷积进行降维。整个网络将多个 Inception 模块堆叠起来，最深的主干达到22层，但通过完全舍弃全连接层（仅在最后使用全局平均池化后接一个线性层分类），总参数量反而只有约500万，远少于 VGGNet。

Inception 模块以多尺度并行计算的方式捕捉不同大小的特征，同时利用 $1\times1$ 卷积实现低成本的特征交互与降维。全局平均池化替代全连接层大幅减少了过拟合风险和参数量。GoogLeNet 在 ILSVRC 2014 上以较低的计算代价取得了与 VGGNet 相当甚至更优的准确率，展现了高效架构设计的巨大潜力。

但多分支结构使网络设计复杂，超参数增多（每个分支的通道数、卷积核大小等），调试和复现需要更多工程经验。后来的经验也表明，过度的多分支不一定总是有益，较新的设计又逐渐回归到更简洁的残差结构。

\paragraph{ResNet}

上述网络在尝试进一步加深时都遇到了瓶颈：层数超过20左右后，训练误差非但不下降，反而升高——这就是著名的“退化”问题。2015年，何恺明等人提出的 ResNet~\cite{he2015deepresiduallearningimage} 以残差连接（residual connection）优雅地解决了这一难题。ResNet 的核心单元是残差块：假设希望让几层网络学习一个目标映射 $\mathcal{H}(\mathbf{x})$，残差块改为学习残差映射 $\mathcal{F}(\mathbf{x}) = \mathcal{H}(\mathbf{x}) - \mathbf{x}$，并在块末尾通过跳跃连接将输入 $\mathbf{x} $ 直接加到输出上，形成 $\mathbf{y} = \mathcal{F}(\mathbf{x}) + \mathbf{x}$。若 $\mathcal{F}(\mathbf{x})$ 全部驱近于零，整个块就退化为恒等映射，因此网络至少不会比浅层版本更差，深层网络得以稳定训练。

ResNet 有多个经典版本，如 ResNet-18、ResNet-34、ResNet-50、ResNet-101、ResNet-152。浅层版本（18/34）使用两个 $3\times3$ 卷积组成的 Basic Block；深层版本（50/101/152）采用 Bottleneck 结构：$1\times1$ 卷积降维、$3\times3$ 卷积、$1\times1$ 卷积升维，以节省计算量。所有网络均以 $7\times7$ 卷积加最大池化开始，随后堆叠残差块，最后以全局平均池化和全连接层分类。ResNet 在 ILSVRC 2015 上以 3.57\% 的 top-5 错误率夺冠，首次超越了人类表现。

残差连接从根本上缓解了梯度消失和退化问题，使训练 100 层以上的网络成为现实。结构仍相对简单规整（尤其是 Basic Block 版本），便于扩展和迁移。ResNet 成为计算机视觉领域的“默认骨干网络”之一，几乎所有的现代架构都或借鉴、或直接在其上扩展。

极深的 ResNet（如152层）训练和推理速度较慢，对算力有较高要求。Bottleneck 结构虽然减少了参数，但依然需要较大的显存来存储中间特征图。此外，残差连接的基本假设是“接近恒等映射有益”，但过量的恒等映射路径也可能限制网络的表达能力，这也是后续 DenseNet 等工作试图改进的方向。


\subsection{目标检测}

目标检测是计算机视觉中的重要领域，旨在识别图像中所有目标对象的位置和类别。与图像分类不同，目标检测不仅需要判断图像中是否存在某类对象，还需要精确地定位这些对象。CNN 在目标检测中的应用主要体现在两个方面：特征提取和候选区域生成。

\paragraph{R-CNN}

R-CNN~\cite{2014Rich} 是将深度学习引入目标检测的开山之作。它的流程分三步走：首先，使用选择性搜索算法在输入图像中生成大约2000个互有重叠的候选区域；然后，将每个候选区域缩放到固定尺寸（$227\times227$），送入一个预训练好的 CNN（如 AlexNet）提取一个固定长度的特征向量；最后，用一组类别特定的线性分类器对这些特征向量进行分类，同时用岭回归器对候选框的坐标进行精修。

在 R-CNN 之前，目标检测完全依赖于手工设计的特征（如 HOG、SIFT）配合滑动窗口，在 VOC 等数据集上性能早已陷入瓶颈。R-CNN 首次证明了 CNN 提取的深度特征能够极大提升检测精度，将 PASCAL VOC 上的平均正确率均值（mAP）一举提升了超过30个百分点。

然而，虽然R-CNN将深度学习引入目标检测任务，但其计算效率极低。每张图像需要独立运行约2000次 CNN 前向传播，这2000个区域大量重叠，造成了严重的重复计算。此外，整个流程分多步独立训练（CNN 微调、SVM 训练、回归器训练），繁琐且难以端到端优化。候选区域必须统一缩放到固定尺寸，也会导致物体的几何失真。

\paragraph{SPPNet}

SPPNet~\cite{He_2014} 的提出正是为了解决 R-CNN 的重复计算和固定尺寸约束这两个问题。它观察到，卷积层本质上是滑窗操作，完全可以在整张图像上只运行一次 CNN，得到全局的特征图，然后根据候选区域在特征图上的投影位置，提取对应的局部特征。更进一步，SPPNet 在最后一个卷积层和全连接层之间插入了一个\textbf{空间金字塔池化}（SPP）层：不管输入区域的特征图尺寸如何，空间金字塔池化都能输出一个固定长度的特征向量。这样既免去了对候选区域逐一缩放裁剪的步骤，也避免了固定输入尺寸的限制。

“一次卷积，反复使用”的策略将 R-CNN 每张图2000次 CNN 前向传播降低为每张图仅1次，整体检测速度提升了两个数量级（约24~102倍，取决于具体配置）。空间金字塔池化又允许网络接受任意尺寸的输入，保留了物体的空间比例信息。

与前代R-CNN相比，SPPNet 的改进主要集中在特征提取阶段，其余部分（SVM分类、回归器、候选区域生成）仍然独立于 CNN 训练，无法实现端到端的整体优化。此外，SPPNet 只在最后一个卷积层上金字塔池化，前面的卷积层在微调时会被冻结，限制了特征学习的能力。

\paragraph{Fast R-CNN}

Fast R-CNN~\cite{7410526} 吸收了 SPPNet“一次卷积”的思想，并在其上实现了质的飞跃。它将空间金字塔池化简化为单尺度的\textbf{感兴趣区域池化}（RoI Pooling）层，将每个候选区域的特征图固定为（例如 $7\times7$）大小。关键是，Fast R-CNN 设计了一个“多任务损失函数”，将类别分类和边界框回归联合训练，并完全舍弃了独立的 SVM 和回归器：网络的末端并行输出 softmax 分类概率和每类边界框偏移。这使得整个网络（卷积层+RoI池化+全连接层）能够以\textbf{端到端}的方式被训练。

Fast R-CNN 成功地将分类和定位统一在单一网络中，训练流程大幅简化，检测精度进一步上升，而训练和推理速度又比 R-CNN 快了数十倍。它彻底证明了目标检测可以像分类一样，纯粹依靠梯度下降来端到端优化。

真正的速度瓶颈此时转移到了外部：选择性搜索生成候选区域本身需要约2秒每张图（GPU上），检测时间仍然远未能达到实时。

\paragraph{Faster R-CNN}

Faster R-CNN~\cite{ren2016fasterrcnnrealtimeobject} 做出了一个大胆而极富洞察力的决定：\textbf{让 CNN 自己来生成候选区域}。它引入了\textbf{区域提议网络}（RPN），这是一个与检测网络共享骨干卷积特征的全卷积小网络。RPN 在特征图的每一个位置处，同时预测多个不同尺度和长宽比的锚框是否为前景，以及需要做的边框修正。RPN 的输出本身就是区域提议，直接送入后续的 RoI 池化和分类/回归网络。自此，整个目标检测流程（区域提议、特征提取、分类、定位）全部在一个神经网络中完成，并可以端到端训练，无需人为干预。

现在，外部候选算法被完全剔除，候选区域生成的时间降至几乎可以忽略不计（约10毫秒）。Faster R-CNN 首次实现了准实时的深度学习目标检测，又在精度上达到了当时的最高水平（PASCAL VOC 和 COCO 上）。

然而，这种手段采用的两阶段检测器虽然精度很高，但在追求极致速度的应用中仍显得复杂。此外，预设的锚框需要手动定义尺度和比例，引入了一定的先验偏差。现在再设计到高精度目标检测，RCNN系列已经隐退江湖，更多的是DETR架构~\cite{carion2020endtoendobjectdetectiontransformers}了。

\paragraph{YOLO}

YOLO（You Only Look Once）~\cite{redmon2016lookonceunifiedrealtime} 彻底抛弃了“区域提议+分类”的两阶段范式，将目标检测重新定义为一个回归问题：网络直接一次性预测图像中所有物体的边界框和类别。具体做法是将输入图像划分为 $S\times S$ 的网格，每个网格负责预测固定数量的边界框及其置信度，以及该网格所含物体的类别概率。整个网络就是一个从输入图像到输出张量的单次前向传播，没有任何中间的区域提议、RoI池化等操作。

YOLO 的速度极快（标准版本可达45 FPS，快速版本超过150 FPS），首次让深度学习检测器真正胜任实时视频流分析。它直接从全局图像出发预测，有效减少了背景误检，对于图像上下文的利用比两阶段方法更好。只不过，YOLO 的网格划分相对粗糙，每个网格只能预测少量边界框，导致它对小物体和密集粘连的物体检测效果较差，空间定位精度也不如基于区域提议的两阶段精细方法。但因为YOLO本身速度实在是太快了，所以时至今日，依然是主流的实时目标检测器之一，迄今已经发展到YOLOv26~\cite{sapkota2026yolo26keyarchitecturalenhancements}，在精度和速度上都有极为显著的提升。

\paragraph{SSD}

SSD（Single Shot MultiBox Detector）~\cite{Liu_2016} 试图结合 Faster R-CNN 的锚框思想和 YOLO 的单阶段回归速度。它在骨干网络（如 VGG-16）的多个不同尺度的特征图上同时进行检测：浅层高分辨率的特征图负责检测小物体，深层低分辨率的特征图负责检测大物体。每个特征图上的每一个位置都预设一组不同长宽比和尺度的锚框，网络直接预测每个锚框的类别置信度和边界框偏移。这既是一次性前向传播（single shot），又借助多尺度特征图缓解了 YOLO 对小物体不敏感的问题。同时通过大量的数据增强（如图像随机裁剪、颜色扭曲）进一步提升鲁棒性。

SSD 在速度与精度的权衡上取得了当时的最佳平衡：速度与 YOLO 相当或更快，而准确率（尤其在多尺度物体的检测上）显著优于 YOLO，甚至能逼近两阶段的 Faster R-CNN。

SSD 需要手动指定不同层锚框的尺度和长宽比，这些超参数对性能影响较大，且不同数据集通常需要不同的配置。此外，用于小物体检测的浅层特征图语义信息相对匮乏，尽管有所改善，但在极端小物体下的性能仍不及精心设计的区域提议方法。

\subsection{人脸识别}

人脸识别也是计算机视觉的一个重要应用领域，涉及从图像或视频中识别和验证个体身份，主要包含两个任务：
\begin{description}
    \item[人脸验证] 判断两张人脸是否属于同一个人，通常用于身份认证。
    \item[人脸识别] 在一个已知人脸数据库中，识别输入人脸属于哪一个已知个体，通常用于考勤、安防等场景。 
\end{description}

在深度学习兴起之前，人脸识别主要依赖于手工设计的局部特征描述子（如 LBP、Gabor）配合高维度的度量学习。这些方法在受控环境下表现尚可，但面对光照、姿态、表情、遮挡等非可控因素时鲁棒性骤降。CNN 的引入彻底改变了这一局面，其核心思路从“手工设计不变特征”转变为“从大量数据中端到端地学习人脸表征”。

\paragraph{DeepFace 与 DeepID：初探深度人脸表征}

2014年的 DeepFace~\cite{2014DeepFace} 是第一个将深度 CNN 应用于大规模人脸识别并获得显著突破的工作。它首先借助三维人脸模型对人脸图像进行对齐（frontalization），将侧脸或带有姿态变化的人脸校正到统一的正面视角；然后将对齐后的 RGB 人脸送入一个包含8层（其中3层为局部连接层）的网络进行训练，最后一层输出一个固定维度的特征向量，并通过 softmax 分类目标驱动学习。DeepFace 在 LFW（Labeled Faces in the Wild）数据集上达到了97.35\%的准确率，首次逼近人类水平。

同期出现的 DeepID~\cite{ouyang2015deepidnetdeformabledeepconvolutional} 系列则走了另一条路：通过训练多个 CNN（每个 CNN 在面部不同区域、不同尺度上学习），将各网络的输出拼接成高维特征向量（DeepID 特征），再训练一个联合贝叶斯分类器进行人脸验证。DeepID2 更进一步引入验证信号（两张人脸是否相同）对分类信号进行辅助，使得学到的特征在类间分离和类内紧凑之间取得更好平衡。DeepID3 则直接将 VGGNet 和 GoogLeNet 的结构移植到人脸特征提取中，进一步将 LFW 准确率推至 99.5\% 以上。

这些工作证明了深度特征能够自动对身份变化、类内差异进行解耦式建模，无需显式的人脸分块或手工特征工程。它们也开启了“在大规模身份标注数据集上预训练、在不同人脸任务上通用”的范式。

\paragraph{FaceNet：三元组损失的范式转变}

2015年，Google 提出的 FaceNet~\cite{Schroff_2015} 将人脸识别带入“直接学习嵌入空间”的阶段，其核心思想至今仍深刻影响着整个领域。FaceNet 并不使用 softmax 分类层，而是通过\textbf{三元组损失}（triplet loss）直接优化特征嵌入：每次输入三个人脸样本——锚样本（anchor）、与锚同人的正样本（positive）、与锚不同人的负样本（negative），目标是让锚与正样本之间的距离小于锚与负样本之间的距离并保持一个余量（margin）。学到的特征经过 L2 归一化后，直接使用欧氏距离进行人脸验证和识别，无需任何额外的分类器。

三元组损失的数学形式为：
\[
\mathcal{L}_{\text{triplet}} = \max\left(0,\; \|f(\mathbf{x}^a) - f(\mathbf{x}^p)\|_2^2 - \|f(\mathbf{x}^a) - f(\mathbf{x}^n)\|_2^2 + \alpha\right),
\]
其中 $f(\cdot)$ 是 CNN 输出的归一化嵌入向量，$\alpha$ 是强制分离的边距。

以往的 softmax 分类方法学到的特征虽然可分，但在类内紧凑性上缺乏显式约束，且特征维度必须与训练集身份数绑定。三元组损失将优化目标直接设定为“同人聚拢、异人远离”，极大提升了特征的泛化能力和验证精度。FaceNet 仅需128维嵌入向量，在 LFW 上即达到99.63\%的准确率，同时在 YouTube Faces DB 等视频人脸识别任务上也登顶当时的榜首。

但是，三元组损失的效果极度依赖于三元组的采样策略。随机采样的大多数三元组已经天然满足距离约束（负样本足够远），对训练几乎没有贡献。因此必须精心挖掘“难三元组”（hard triplets），即 $\|f(\mathbf{x}^a)-f(\mathbf{x}^p)\|_2 \approx \|f(\mathbf{x}^a)-f(\mathbf{x}^n)\|_2$ 的那些样本。这一挖掘过程计算开销巨大，直接影响收敛速度和最终性能。

\paragraph{损失函数的演进：从 Softmax 变体到 ArcFace}

FaceNet 之后，研究者们意识到三元组损失的训练复杂度较高，转而重新审视基于分类的损失函数是否可以兼具“类间分离”和“类内紧凑”两种特性。一系列改进的 softmax 变体应运而生。

先从基础 softmax 损失的形式入手：
\[
\mathcal{L}_{\text{softmax}} = -\frac{1}{N}\sum_{i=1}^N \log\frac{e^{\mathbf{W}_{y_i}^\top \mathbf{x}_i + b_{y_i}}}{\sum_{j=1}^C e^{\mathbf{W}_j^\top \mathbf{x}_i + b_j}},
\]
其中 $\mathbf{x}_i$ 为样本 $i$ 的特征向量，$\mathbf{W}_j$ 为第 $j$ 类的权重向量，$C$ 为类别总数。这个目标的本质是让 $\mathbf{W}_{y_i}^\top \mathbf{x}_i$（正确类的内积，可视为余弦相似度乘以模长）尽可能大。

2017年的工作~\cite{Wang_2017,wang2018cosfacelargemargincosine} 相继指出，偏置项和权重模长可以通过归一化消除，使得优化目标归结为纯粹的余弦相似度，避免模型依赖“增加某些类别权重的模长”来降低损失。将 $\|\mathbf{W}_j\|=1$、$b_j=0$ 固定后，决策边界仅由角度决定。

在归一化的基础上，为了增强类内紧凑性，可以引入角度边距（angular margin）。L-Softmax~\cite{liu2017largemarginsoftmaxlossconvolutional} 通过放缩因子收紧正确类别的分类边界；SphereFace~\cite{liu2018spherefacedeephypersphereembedding} 进一步修改为 $\mathbf{W}_{y_i}^\top \mathbf{x}_i = \|\mathbf{x}_i\|\cos(m\theta_{y_i})$，直接将角度的倍数作为边距，不同类别的特征向量被推向角度空间的更远处。

CosFace~\cite{wang2018cosfacelargemargincosine} 使用余弦边距（cosine margin），即 $\mathbf{W}_{y_i}^\top \mathbf{x}_i = \cos\theta_{y_i} - m$；ArcFace~\cite{Deng_2022} 则使用加性角度边距（additive angular margin），即 $\mathbf{W}_{y_i}^\top \mathbf{x}_i = \cos(\theta_{y_i} + m)$。两者的思想极其相似：在正确类与特征向量的余弦相似度上直接减去一个边距常数（或加上角度常数），使得分类决策边界更加严苛，迫使同类特征向中心聚拢。ArcFace 的训练稳定、收敛快，在各类基准上取得了当时的最佳性能。

这些改进的 softmax 变体综合了 softmax 的类间分离能力和度量学习（如三元组损失）的类内紧凑性要求，同时完全避免了困难样本挖掘的工程负担。它们使得大规模人脸识别的训练流程标准化、可复现，学到的嵌入向量可以直接用余弦相似度进行比对，简单高效。ArcFace 自此成为人脸识别领域事实上的标准损失函数。

近年来，人脸识别的核心演进逐渐从模型结构转向\textbf{数据规模}与\textbf{训练基础设施}。借助数千万乃至数亿张身份标注图像、复杂的自动清洗与去重流水线、分布式并行训练，业界已将人脸识别系统在百万人级别开放集上的识别精度推动到了实用化高度。主流框架（如基于 ResNet 的嵌入网络 + ArcFace 损失 + 大余量微调）已经成为各类人脸识别产品的标配。
